%--------------------------------------------------------------------%
% BAGIAN I. DOKUMEN SEMINAR PROPOSAL
% Wajib pada seminar proposal. Pada sidang akhir bagian ini digantikan
% oleh Bab I sampai VIII yang final (nonaktifkan \include di thesis.tex).
%--------------------------------------------------------------------%
\halamanbagian{BAGIAN I. DOKUMEN SEMINAR PROPOSAL}{Wajib pada seminar proposal, sebelum pelaksanaan. Pada sidang akhir bagian ini digantikan oleh Bab I sampai VIII yang final.}

%--------------------------------------------------------------------%
\judulhalaman{P.1 PROFIL DAN LEGALITAS MITRA}
\petunjuk{Mitra harus berbadan hukum (Akte Pendirian, SK Lembaga, NPWP Badan Usaha). Surat Pernyataan Kesediaan Mitra ditandatangani \textit{Top Management} mitra sebelum seminar proposal dan dilampirkan pada Lampiran 2 (Pedoman TA butir 3.3.3). Uraikan secara naratif: nama lembaga, bentuk badan hukum, alamat, \textit{Top Management} (nama dan jabatan), nomor Akte Pendirian dan SK Lembaga, NPWP badan usaha, dan kontak mitra. \tahap{\wajib{WAJIB}}{\tidakperlu{Tidak diperlukan}}}
\tulisuraian

%--------------------------------------------------------------------%
\judulhalaman{P.2 TIM DAN PEMBAGIAN KERJA}
\petunjuk{Kelompok 3 sampai 4 mahasiswa; seminar dilaksanakan bersamaan untuk seluruh anggota. Uraikan nama, NIM, peran, tugas, dan judul Laporan TA individu tiap anggota; judul laporan individu wajib sudah ditetapkan saat seminar proposal. Lampirkan Surat Pernyataan Kesanggupan Pembagian Kerja Tim beserta \textit{timeline} pada Lampiran 3 (Pedoman TA butir 3.3.2 dan 5.1.3). \tahap{\wajib{WAJIB}}{\tidakperlu{Tidak diperlukan}}}
\tulisuraian

%--------------------------------------------------------------------%
\judulhalaman{P.3 RENCANA PELAKSANAAN}
\petunjuk{Rancang pekerjaan agar selesai dalam 3 semester aktif sejak seminar proposal (Pedoman TA butir 3.3.2). \tahap{\wajib{WAJIB}}{\tidakperlu{Tidak diperlukan}}}

\subbabP{P.3.1\quad Jadwal}
\petunjuk{Uraikan kegiatan, waktu mulai dan selesai, serta penanggung jawab untuk tahap analisis kebutuhan, perancangan solusi, implementasi prototipe, pengujian dan evaluasi, evaluasi dan serah terima mitra, serta sidang akhir.}
\tulisuraian

\subbabP{P.3.2\quad Anggaran dan Sumber Daya}
\petunjuk{Anggaran dan sumber daya yang dibutuhkan beserta penanggungnya.}
\tulisuraian

\subbabP{P.3.3\quad Rencana Pengujian Menuju TKT 3}
\petunjuk{Uji yang akan dilakukan, kriteria lulus, dan cara membuktikan TKT minimal 3 tercapai (Pedoman TA butir 3.3.2).}
\tulisuraian

\subbabP{P.3.4\quad Risiko dan Rencana Cadangan}
\petunjuk{Risiko utama dan rencana cadangannya, misalnya anggota berhalangan atau mitra tidak dapat hadir. \tahap{\opsional{Opsional}}{\tidakperlu{Tidak diperlukan}}}
\tulisuraian
